% part: intuitionistic-logic
% chapter: propositions-as-types
% section: types

\documentclass[../../../include/open-logic-section]{subfiles}

\begin{document}

\olfileid{int}{pty}{typ}

\olsection{نوع‌ها}

ترم برهانی مانند~$(MN)$ به‌تنهایی برهانی را که ترم آن است به‌طور یکتا
تعیین نمی‌کند. علت این است که متغیرهای آزاد، یعنی ترم‌های برهان اصول،
!!{formula}s موجود در اصول متناظر را مشخص نمی‌کنند. اما اگر این
اطلاعات را در بافتی چون~$\Gamma$ تعیین کنیم، برهان، در صورت وجود،
یکتا تعیین می‌شود.

\begin{defn}
یک \emph{بافت} برای ترم برهان~$N$ مجموعه‌ای چون~$\Gamma$ است که برای
هر متغیر آزاد~$x$ در~$N$ دقیقاً یک جفت~$x:!A$ در بر دارد.
\end{defn}

توجه کنید تعریف لازم نمی‌داند~$\Gamma$ \emph{فقط} برای متغیرهای آزاد
~$N$ جفت~$x:!A$ داشته باشد. بنابراین، برای نمونه،
$\Gamma=\{x:!A,y:!B,z:!C\}$ بافتی برای~$\pair{x}{y}$ است.

\begin{defn}\ollabel{defn:type}
  فرض کنید~$\Gamma$ بافتی برای~$N$ باشد. \emph{نوع}~$N$ نسبت
  به~$\Gamma$ با بندهای استقرایی زیر تعریف می‌شود:
  \begin{enumerate}
  \item نوع~$x$ برابر~$!A$ است اگر و تنها اگر~$x:!A \in \Gamma$.
  \item اگر~$N$ نسبت به~$x:!A,\Gamma$ نوع~$!B$ داشته باشد، نوع
    $\lambd[\typeof{x}{!A}][N]$ برابر~$!A \lif !B$ است.
  \item اگر~$N$ نوع~$!A \lif !B$ و~$M$ نوع~$!A$ داشته باشد، آنگاه
    $(NM)$ نوع~$!B$ دارد.
  \item اگر~$N$ نوع~$!A$ و~$M$ نوع~$!B$ داشته باشد، آنگاه
    $\pair{N}{M}$ نوع~$!A \land !B$ دارد.
  \item اگر~$N$ نوع~$!A_1 \land !A_2$ داشته باشد، آنگاه
    $\proj{i}{N}$ نوع~$!A_i$ دارد.
  \item اگر~$N$ نوع~$!A$ داشته باشد، آنگاه
    $\inj[!B]{i}{N}$ برای~$i=1$ نوع~$!A \lor !B$ و برای~$i=2$
    نوع~$!B \lor !A$ دارد.
  \item اگر~$M$ نوع~$!A \lor !B$، و~$N_1$ نسبت به
    $x_1:!A,\Gamma$ نوع~$!C$ و~$N_2$ نسبت به
    $x_2:!B,\Gamma$ نوع~$!C$ داشته باشند، آنگاه
    $\dcase{M}{x_1}{N_1}{x_2}{N_2}$ نوع~$!C$ دارد.
  \item اگر~$N$ نوع~$\lfalse$ داشته باشد، آنگاه
    $\abort{!A}{N}$ نوع~$!A$ دارد.
  \end{enumerate}
\end{defn}

\begin{prop}
نسبت به هر بافت~$\Gamma$ برای~$N$، ترم برهان~$N$ حداکثر یک نوع دارد؛
یعنی اگر نوعی برای آن وجود داشته باشد، آن نوع یکتا است.
\end{prop}

\begin{proof}
  بر~$N$ استقرا می‌کنیم.
\end{proof}

اگر نوع~$N$ نسبت به~$\Gamma$ برابر~$!A$ باشد، می‌نویسیم
$\Gamma \Sequent N:!A$. بندهای~\olref{defn:type} باید آشنا به نظر
برسند: آن‌ها دقیقاً قواعد~\Log{tN2i} هستند، با یک تفاوت کوچک؛ بافت
~$\Gamma$ در سراسر استنتاج یکسان است. می‌توانیم تعریف را به‌صورت حساب
~\Log{tN3i} با قواعد مندرج در~\olref{tab:tN3ip} صورت‌بندی کنیم.

\subfile{rules-tN3}

ترم‌های برهان، ترم‌های گونه‌ای از حساب~$\lambd$ به نام
\emph{حساب~$\lambd$ نوع‌دار با ضرب‌ها، جمع‌ها و نوع تهی} هستند. حساب
~$\lambd$ کلاسیک تنها ترم‌هایی دارد که از متغیرها،
\emph{انتزاع‌های~$\lambd$} مانند~$\lambd[x][N]$ و اعمال~$(MN)$ ساخته
شده‌اند؛ و حاشیه‌نویسی~$!A$ در انتزاع نوع‌دار
$\lambd[\typeof{x}{!A}][N]$ را ندارد. حساب~$\lambd$ مدلی تورینگ‌کامل
از محاسبه است؛ یعنی هر تابع جزئی محاسبه‌پذیر را می‌توان با ترمی در آن
بازنمایی کرد. گونه‌ای از حساب~$\lambd$ که ترم‌های برهان و قواعد نوع‌دهی
می‌سازند مدلی محدودتر از محاسبه است، اما اندیشه‌های اصلی همان‌اند.
انتساب نوع آن را محدود می‌کند: تنها ترم‌های \emph{درست‌نوع‌دار}، یعنی
ترم‌های برهانی که نسبت به بافتی چون~$\Gamma$ نوع دارند، ترم‌های مجازند.
برای نمونه،~$\lambd[x][(xx)]$ در حساب~$\lambd$ کلاسیک بی‌نوع مجاز است،
اما~$\lambd[\typeof{x}{!A}][(xx)]$ برای هیچ~$!A$ درست‌نوع‌دار نیست.
زیرا قواعد نوع‌دهی برای نوع‌داشتن~$(xx)$ ایجاب می‌کنند نخستین~$x$ نوع
$!A \lif !B$ و دومین~$x$ نوع~$!A$ داشته باشد؛ این ناممکن است، چون
$!A$ یک زیر-!!{formula} سرهٔ~$!A \lif !B$ است.

به‌طور شهودی می‌توان نوع‌ها را گونهٔ چیزی دانست که ترمی از آن نوع
می‌نامد. پس ترمی از نوع~$!A \land !B$ جفتی را برمی‌گزیند که مؤلفهٔ
نخستش نوع~$!A$ و مؤلفهٔ دومش نوع~$!B$ دارد؛ یعنی !!a{element} از
ضرب~$!A \times !B$. ترمی از نوع~$!A \lif !B$ تابعی از چیزهای نوع
~$!A$ به چیزهای نوع~$!B$ است. ترمی از نوع~$!A \lor !B$ یا چیزی از
نوع~$!A$ است یا چیزی از نوع~$!B$، همراه با اطلاعاتی که مشخص می‌کند
کدام‌یک است. این همانند اجتماع مجزای دو مجموعهٔ~$!A$ و~$!B$ است که
چنین تعریف می‌شود:
$\{\tuple{i,x}:i=1\text{ یا }2,\ x\in !A\text{ اگر }i=1,\ x\in !B
\text{ اگر }i=2\}$. نوع~$\lfalse$ مجموعهٔ تهی است؛ یعنی ترمی با این
نوع نمی‌تواند چیزی را بنامد. چون استنتاج طبیعی، از جمله~\Log{tN3i}،
سازگار است، هیچ !!{proof}ی از~$\lfalse$ بدون فرض‌های باز وجود ندارد؛
پس هیچ ترم برهان بسته‌ای، یعنی ترمی بدون متغیر آزاد، نوع~$\lfalse$
ندارد.

عملگرهای حساب~$\lambd$ نوع‌دار ترم‌هایی می‌سازند که، به‌طور شهودی،
چیزهایی از گونهٔ درست را می‌نامند، مشروط بر آنکه ترم‌های ورودی‌شان
چیزهای معینی را بنامند. برای نمونه، «اگر~$N_1:!A$ و~$N_2:!B$، آنگاه
$\pair{N_1}{N_2}:!A\land !B$» می‌گوید اگر~$N_1$ چیزی از نوع~$!A$ و
~$N_2$ چیزی از نوع~$!B$ را بنامد، آنگاه~$\pair{N_1}{N_2}$ چیزی از
نوع~$!A \land !B$ را می‌نامد.

\end{document}
